% from ega1-8-fr.tex
\oldpage[I]{169}
tel que $\Gamma(U,\mathscr{O}_X)=\bigcap_{x\in U}\mathscr{O}_x$ pour tout
ouvert $U\subset X$, avec une définition évidente des homomorphismes de
restriction. Nous laissons au lecteur le soin de vérifier qu'on obtient bien
ainsi un schéma intègre, dont les anneaux locaux sont les éléments de $X'$~;
nous n'utiliserons pas ce résultat par la suite.
\end{env}

% (continued in ega1-9-fr.tex, before its first page marker)
\section{Compléments sur les faisceaux quasi-cohérents}
\label{section:I.9-fr}

\subsection{Produit tensoriel de faisceaux quasi-cohérents.}
\label{subsection:I.9.1-fr}

\begin{proposition}[9.1.1]
\label{I.9.1.1-fr}
Soit $X$ un préschéma (resp. un préschéma localement noethérien). Soient
$\mathscr{F}$ et $\mathscr{G}$ deux $\mathscr{O}_X$-Modules quasi-cohérents
(resp. cohérents)~; alors
$\mathscr{F}\otimes_{\mathscr{O}_X}\mathscr{G}$ est quasi-cohérent (resp.
cohérent) et de type fini si $\mathscr{F}$ et $\mathscr{G}$ sont de type
fini. Si $\mathscr{F}$ admet une présentation finie et si $\mathscr{G}$ est
quasi-cohérent (resp. cohérent), alors
$\mathscr{H}\!om_{\mathscr{O}_X}(\mathscr{F},\mathscr{G})$ est quasi-cohérent
(resp. cohérent).
\end{proposition}

La question étant locale, on peut supposer $X$ affine (resp. affine
noethérien)~; en outre, si $\mathscr{F}$ est cohérent, on peut supposer qu'il
est le conoyau d'un homomorphisme
$\mathscr{O}_X^m\to\mathscr{O}_X^n$. Les assertions relatives aux faisceaux
quasi-cohérents résultent alors de
(\hyperref[I.1.3.12-fr]{1.3.12}) et
(\hyperref[I.1.3.9-fr]{1.3.9})~; les assertions relatives aux faisceaux
cohérents résultent de (\hyperref[I.1.5.1-fr]{1.5.1}) et du fait que, si $M$
et $N$ sont des modules de type fini sur un anneau noethérien $A$,
$M\otimes_A N$ et $\operatorname{Hom}_A(M,N)$ sont des $A$-modules de type
fini.

\begin{definition}[9.1.2]
\label{I.9.1.2-fr}
Soient $X$, $Y$ deux $S$-préschémas, $p$, $q$ les projections de
$X\times_S Y$, $\mathscr{F}$ (resp. $\mathscr{G}$) un
$\mathscr{O}_X$-Module (resp. un $\mathscr{O}_Y$-Module) quasi-cohérent. On
appelle produit tensoriel de $\mathscr{F}$ et $\mathscr{G}$ sur
$\mathscr{O}_S$ (ou sur $S$) et on note
$\mathscr{F}\otimes_{\mathscr{O}_S}\mathscr{G}$ (ou
$\mathscr{F}\otimes_S\mathscr{G}$) le produit tensoriel
$p^*(\mathscr{F})\otimes_{\mathscr{O}_{X\times_S Y}}q^*(\mathscr{G})$ sur le
préschéma $X\times_S Y$.
\end{definition}

Si $X_i$ ($1\leq i\leq n$) sont des $S$-préschémas, $\mathscr{F}_i$ un
$\mathscr{O}_{X_i}$-Module quasi-cohérent ($1\leq i\leq n$), on définit de
la même manière le produit tensoriel
$\mathscr{F}_1\otimes_S\mathscr{F}_2\otimes_S\cdots\otimes_S\mathscr{F}_n$
sur le préschéma
$Z=X_1\times_S X_2\times_S\cdots\times_S X_n$~; c'est un
$\mathscr{O}_Z$-module \emph{quasi-cohérent} en vertu de
(\hyperref[I.9.1.1-fr]{9.1.1}) et de
(\hyperref[0.5.1.4-fr]{0, 5.1.4})~; il est \emph{cohérent} si les
$\mathscr{F}_i$ le sont et si $Z$ est \emph{localement noethérien} en vertu
de (\hyperref[I.9.1.1-fr]{9.1.1}),
(\hyperref[0.5.3.11-fr]{0, 5.3.11}) et
(\hyperref[I.6.1.1-fr]{6.1.1}).

On notera que si on prend $X=Y=S$, la définition
(\hyperref[I.9.1.2-fr]{9.1.2}) redonne le produit tensoriel de
$\mathscr{O}_S$-Modules. En outre, comme
$q^*(\mathscr{O}_Y)=\mathscr{O}_{X\times_S Y}$
(\hyperref[0.4.3.4-fr]{0, 4.3.4}), le produit
$\mathscr{F}\otimes_S\mathscr{O}_Y$ s'identifie canoniquement à
$p^*(\mathscr{F})$, et de même
$\mathscr{O}_X\otimes_S\mathscr{G}$ s'identifie canoniquement à
$q^*(\mathscr{G})$. Plus particulièrement, si on prend $Y=S$ et qu'on
désigne par $f$ le morphisme structural $X\to Y$, on a
$\mathscr{O}_X\otimes_Y\mathscr{G}=f^*(\mathscr{G})$~: le produit tensoriel
ordinaire et l'image réciproque apparaissent donc comme des cas particuliers
du produit tensoriel général.

La déf. (\hyperref[I.9.1.2-fr]{9.1.2}) entraîne immédiatement que, pour $X$
et $Y$ fixés, $\mathscr{F}\otimes_S\mathscr{G}$ est un
\emph{bifoncteur covariant additif et exact à droite} en $\mathscr{F}$ et
$\mathscr{G}$.

\begin{proposition}[9.1.3]
\label{I.9.1.3-fr}
Soient $S$, $X$, $Y$ trois schémas affines d'anneaux respectifs, $A$, $B$,
$C$, $B$ et $C$ étant des $A$-algèbres. Soit $M$ (resp. $N$) un $B$-module
(resp. $C$-module), $\mathscr{F}=\widetilde{M}$ (resp.
$\mathscr{G}=\widetilde{N}$) le faisceau quasi-cohérent associé~; alors
$\mathscr{F}\otimes_S\mathscr{G}$ est canoniquement isomorphe au faisceau
associé au $(B\otimes_A C)$-module $M\otimes_A N$.
\end{proposition}

\begin{proof}
